% =====================================================================
%  Préambule — mise en forme, packages, commandes maison.
%  Rien de spécifique au contenu ici : le texte vit dans chapters/.
% =====================================================================

% ---------- Langue et écriture ----------
%
%  COMPILER AVEC XeLaTeX  (xelatex main.tex), pas pdflatex.
%
%  Le modèle de l'institut impose un résumé en arabe sur la dernière page. pdfLaTeX ne sait pas
%  composer l'arabe : il ne gère ni le sens droite-à-gauche ni les ligatures contextuelles. On passe
%  donc par XeTeX, fontspec et polyglossia, qui lisent directement les polices système. C'est aussi
%  pourquoi babel, inputenc et fontenc ont disparu : ils sont incompatibles avec polyglossia et
%  inutiles sous XeTeX, dont l'entrée est nativement UTF-8.
\usepackage{fontspec}
\usepackage{polyglossia}
\setmainlanguage{french}
\setotherlanguage{arabic}
% Times New Roman, corps 12 : imposé par le guide de rédaction de l'institut (règle 1).
\setmainfont{Times New Roman}
\newfontfamily\arabicfont[Script=Arabic]{Amiri}  % naskh académique, fourni par MiKTeX
\usepackage{csquotes}          % requis par biblatex pour les guillemets

% ---------- Typographie ----------
\usepackage{microtype}
\usepackage{setspace}
\onehalfspacing                % interligne 1,5 — attendu par la plupart des jurys
\usepackage{parskip}

% ---------- Géométrie ----------
% Largeur et hauteur de texte imposées par le guide de rédaction (règle 3) : 160 × 250 mm sur A4
% (210 × 297 mm). Ne pas fixer les marges directement : sans elles, geometry centre le bloc de
% texte pour obtenir cette largeur et cette hauteur exactement, ce qui est le résultat demandé.
% bindingoffset décale le bloc vers la droite pour la reliure collée, sans toucher à sa largeur :
% le pli en avale environ 5 mm, qui seraient sinon pris sur la marge de gauche.
% headheight : fancyhdr reclame 14.5pt et n'en recoit que 12 par defaut, ce qu'il signale a chaque
% page. L'option est passee a geometry plutot que fixee apres coup : geometry ajuste alors la marge
% haute et laisse textheight a 250mm exactement, comme la regle 3 l'impose. La fixer par
% \setlength apres le chargement ferait deborder le bloc de 2.5pt vers le bas.
\usepackage[a4paper,textwidth=160mm,textheight=250mm,bindingoffset=5mm,headheight=14.5pt]{geometry}

% ---------- Figures et tableaux ----------
\usepackage{graphicx}
% Les diagrammes sont ranges par sprint depuis leur reorganisation : chaque dossier est
% declare ici, sinon \includegraphics ne trouve rien a la racine de drawio/.
\graphicspath{{figures/}{figures/logos/}{../diagrams/drawio/}{../diagrams/drawio/Sprint0/}
  {../diagrams/drawio/Sprint1/}{../diagrams/drawio/Sprint2/}{../diagrams/drawio/Sprint3/}
  {../diagrams/drawio/Sprint4/}{../diagrams/drawio/Sprint5/}{../diagrams/drawio/Architecture/}
  {../diagrams/drawio/chap9/}{../diagrams/svg/}{../docs/assets/uml/}}
% Les diagrammes UML sont reduits pour tenir dans le bloc de texte, et leurs libelles descendent
% alors sous la taille lisible. Le bloc fait 160 mm sur une page de 210 : une figure peut occuper
% 190 mm sans toucher le bord, soit 19 % de texte en plus, sans rien changer au diagramme. Le
% \makebox garde la figure centree sur le bloc et empeche LaTeX de signaler un debordement.
\newcommand{\diagramme}[2][0.45]{%
  \makebox[\textwidth][c]{\includegraphics[width=190mm,height=#1\textheight,keepaspectratio]{#2}}}
% Les deux diagrammes de sequence les plus larges restaient illisibles en portrait : reduits pour
% tenir dans 190 mm, leurs libelles descendaient sous 4 pt. Ils occupent desormais une page entiere
% en paysage, ou la largeur utile passe a 280 mm. pdflscape fait pivoter la page dans le lecteur,
% de sorte que la figure se lit sans tourner le document imprime.
\usepackage{pdflscape}
% La hauteur est la dimension contraignante pour ces deux diagrammes : en paysage, le bloc de texte
% n'en offre que 160 mm alors que la page tournee en fait 210, marges comprises. On lui en rend 15,
% ce qui laisse 10 mm sous la legende : la figure passe de 132 mm en portrait a 163, soit 24 % de
% plus. Au-dela, la legende sort de la feuille sans que LaTeX ne signale rien.
\newcommand{\diagrammepaysage}[3]{%
  \begin{landscape}
  \enlargethispage{15mm}%
  \begin{figure}[H]\centering
  \makebox[\textwidth][c]{\includegraphics[width=270mm,height=163mm,keepaspectratio]{#1}}
  \caption{#2}\label{#3}
  \end{figure}
  \end{landscape}}
\usepackage{float}
% Les figures en [H] laissaient de grands blancs : quand l'image ne tenait plus dans la place
% restante, LaTeX abandonnait le bas de la page. En [htbp] elles flottent vers le haut ou le bas de
% la page, et le texte comble. placeins/section les empeche de deriver au-dela de leur section, de
% sorte qu'une figure reste toujours dans la section qui l'annonce.
\usepackage[section]{placeins}
% placeins ne pose sa barriere qu'aux \section. Sans cette extension aux \subsection, une figure
% annoncee dans une sous-section peut deriver jusqu'a la fin de la section : le diagramme de cas
% d'utilisation du sprint 3 arrivait quatre pages apres la phrase qui l'annonce, une fois passes
% les trois tableaux de cas detailles.
\makeatletter
\AtBeginDocument{%
  \expandafter\renewcommand\expandafter\subsection\expandafter{%
    \expandafter\@fb@secFB\subsection}%
}
\makeatother
% Seuils de flottement. Par defaut LaTeX refuse qu'une figure occupe plus de 70 % du haut d'une page
% et exige 20 % de texte, ce qui reporte les grandes images et laisse des demi-pages blanches. On lui
% permet d'aller jusqu'a 92 % en haut, de n'exiger que 5 % de texte, et de remplir une page de
% flottants des 60 %, au lieu de 50.
\renewcommand{\topfraction}{0.92}
\renewcommand{\bottomfraction}{0.85}
\renewcommand{\textfraction}{0.05}
\renewcommand{\floatpagefraction}{0.60}
\setcounter{topnumber}{3}
\setcounter{bottomnumber}{2}
\setcounter{totalnumber}{5}
% Espacement autour des flottants place en [H] : la valeur par defaut de LaTeX est genereuse et
% suffisait, sur une page deja chargee, a faire basculer la figure suivante a la page d'apres en
% laissant un grand blanc. La resserrer fait tenir un couple texte + capture de plus par page.
\setlength{\intextsep}{10pt plus 2pt minus 3pt}
\setlength{\textfloatsep}{12pt plus 2pt minus 4pt}
\usepackage{caption}
\usepackage{subcaption}
\usepackage{booktabs}
\usepackage{longtable}
\usepackage{tabularx}
\usepackage{multirow}
\usepackage{array}
\usepackage[table]{xcolor}

% Colonnes de largeur fixe, alignées à gauche / centrées / à droite.
\newcolumntype{L}[1]{>{\raggedright\arraybackslash}p{#1}}
\newcolumntype{C}[1]{>{\centering\arraybackslash}p{#1}}
\newcolumntype{R}[1]{>{\raggedleft\arraybackslash}p{#1}}

% ---------- Mise en page libre (couverture, page de résumé) ----------
\usepackage{tikz}
\usetikzlibrary{calc}

% ---------- Couleurs ----------
% Les deux premières sont relevées sur la charte ISIMS (bande de couverture) : la couverture doit
% rester conforme au modèle de l'institut, ce n'est pas un choix esthétique.
\definecolor{isimsblue}{HTML}{1F6CA0}
\definecolor{isimsyellow}{HTML}{D9DE28}
\definecolor{primary}{HTML}{1B5E20}
\definecolor{accent}{HTML}{00897B}
\definecolor{codebg}{HTML}{F5F5F5}

% ---------- En-têtes ----------
\usepackage{fancyhdr}
\pagestyle{fancy}
\fancyhf{}
\fancyhead[L]{\nouppercase{\leftmark}}
\fancyhead[R]{\thepage}
\renewcommand{\headrulewidth}{0.4pt}
\fancypagestyle{plain}{\fancyhf{}\fancyfoot[C]{\thepage}\renewcommand{\headrulewidth}{0pt}}

% ---------- Titres de chapitre ----------
% Motif carré noir + numéro blanc, titre à côté, filet en dessous : observé identique dans deux
% mémoires ISIMS réels (couverture de ce document et le mémoire InvisiThreat pris comme modèle),
% sur les titres de chapitre ET les titres de section non numérotés (Dédicace, Remerciements, Table
% des matières...). C'est la convention de l'institut, pas un choix esthétique — d'où l'abandon du
% vert utilisé dans une version précédente : rien dans le modèle observé n'est en couleur.
%
% titlesec ne permet pas facilement de mettre une étiquette et un titre côte à côte à la même
% hauteur (son mode [display] les empile verticalement) ; on redéfinit donc directement les
% macros internes de la classe report, qui donnent un contrôle total sur la mise en page.
\newcommand{\chapitrebadge}[1]{%
  \tikz[baseline=-0.5ex]{\node[fill=black,text=white,font=\bfseries\Large,
    minimum width=1.2cm,minimum height=1.2cm,inner sep=0pt]{#1};}%
}
\makeatletter
% Chapitre numéroté (corps du mémoire).
\def\@makechapterhead#1{%
  \vspace*{10\p@}%
  \parindent\z@\raggedright
  \begin{minipage}[b]{2.1cm}
    % \chaptertitlename n'existe pas en LaTeX standard — c'est un alias fourni par titlesec, qui
    % n'est plus chargé ici. \@chapapp est la vraie macro du noyau (report.cls) : elle vaut
    % « Chapitre » ou « Annexe » selon qu'on est dans \appendix ou non, et suit la langue active.
    % Largeur 2,1cm : à 1,5cm, « CHAPITRE » en dépassait de 0,5cm (overfull hbox).
    \normalfont\small\bfseries\MakeUppercase{\@chapapp}\\[6pt]
    \chapitrebadge{\thechapter}
  \end{minipage}%
  \hfill
  \begin{minipage}[b]{0.70\textwidth}
    \raggedright\normalfont\Large\bfseries #1
  \end{minipage}\par\nobreak
  \vskip 10\p@
  \noindent\rule{\textwidth}{1pt}
  \vskip 24\p@
}
% Chapitre non numéroté (\chapter*{...} : Dédicace, Remerciements, Annexes...) : même filet, mais
% un carré plein sans numéro à la place du badge, comme dans le modèle.
\def\@makeschapterhead#1{%
  \vspace*{10\p@}%
  \parindent\z@\raggedright
  \tikz{\node[fill=black,minimum width=0.4cm,minimum height=1.1cm,inner sep=0pt]{};}%
  \hspace{10pt}%
  {\normalfont\Large\bfseries #1}\par\nobreak
  \vskip 10\p@
  \noindent\rule{\textwidth}{1pt}
  \vskip 24\p@
}
\makeatother

% ---------- Code source ----------
\usepackage{listings}
\lstset{
  basicstyle=\ttfamily\footnotesize,
  backgroundcolor=\color{codebg},
  keywordstyle=\color{primary}\bfseries,
  commentstyle=\color{gray}\itshape,
  stringstyle=\color{accent},
  numbers=left, numberstyle=\tiny\color{gray}, numbersep=8pt,
  breaklines=true, showstringspaces=false,
  frame=single, rulecolor=\color{gray!40},
  captionpos=b,
  % Les extraits contiennent des accents (commentaires en français).
  extendedchars=true, literate={é}{{\'e}}1 {è}{{\`e}}1 {à}{{\`a}}1 {ê}{{\^e}}1 {ç}{{\c c}}1 {û}{{\^u}}1 {ô}{{\^o}}1 {î}{{\^\i}}1
}

% ---------- Éviter un pied de page orphelin ----------
\usepackage{needspace}

% ---------- Listes ----------
\usepackage{enumitem}
\setlist{itemsep=2pt,topsep=4pt}

% ---------- Liens ----------
\usepackage[hidelinks]{hyperref}
\hypersetup{
  pdftitle={ASM Track — Rapport de Projet de Fin d'Études},
  pdfauthor={Mohamed Aziz Hadj Kacem},
  bookmarksnumbered=true
}

% ---------- Bibliographie ----------
% Le guide de rédaction (règle 4) impose un corps 10 et un style par abréviation entre crochets :
% [ROU 99] pour Paul Rouet, 1999, soit trois lettres du nom, une espace, deux chiffres de l'année.
% \bibfont est le point d'entrée officiel de biblatex pour la taille de police de la liste.
% maxalphanames=1 : l'etiquette ne retient que le premier auteur, sans quoi biblatex accole trois
% lettres de chacun (SCHSUT pour Schwaber et Sutherland). labelalphaothers vide supprime le « + »
% que le style ajoute des qu'une liste d'auteurs est tronquee (SAN+96).
\usepackage[backend=biber,style=alphabetic,sorting=none,
            maxalphanames=1,minalphanames=1]{biblatex}
\renewcommand*{\labelalphaothers}{}
\renewcommand*{\bibfont}{\small}          % ≈ corps 10 sur un document en corps 12
\addbibresource{bibliography/refs.bib}

% Le style « alphabetic » n'applique la règle des trois lettres qu'à un auteur unique : dès qu'il y
% en a deux il concatène leurs initiales (SS20 pour Schwaber et Sutherland), et au-delà de trois il
% ajoute un « + » (San+96). Aucune des deux formes n'est celle que le guide demande. On redéfinit
% donc l'étiquette pour qu'elle prenne toujours les trois premières lettres du premier nom, en
% capitales, quel que soit le nombre d'auteurs.
\DeclareLabelalphaTemplate{
  \labelelement{
    \field[strwidth=3,strside=left,uppercase=true]{labelname}
    \field[strwidth=3,strside=left,uppercase=true]{label}
    \field[strwidth=3,strside=left,uppercase=true]{shorthand}
    \field[strwidth=3,strside=left,uppercase=true]{title}
  }
  % \space est absorbe par biber ; \addnbspace survit jusqu'au .bbl et produit l'espace insecable
  % qui separe les lettres des chiffres, comme dans le [ROU 99] du guide.
  \labelelement{
    \literal{\addnbspace}
  }
  \labelelement{
    \field[strwidth=2,strside=right]{year}
  }
}
% Sans cela, biblatex ajoute un suffixe « a », « b »… aux etiquettes homonymes, ce que le format
% impose de ne pas faire tant qu'aucune collision reelle n'existe dans cette bibliographie.
\DeclareLabelalphaNameTemplate{
  \namepart{family}
}

% ---------- Raccourcis maison ----------
% Un logo de technologie ou d'outil. Les fichiers melent des icones carrees et des logotypes en
% bandeau : borner la seule hauteur laissait les seconds deborder la marge, d'ou le plafond sur les
% deux dimensions avec conservation des proportions.
\newcommand{\logo}[1]{\includegraphics[height=1cm,width=2.9cm,keepaspectratio]{#1}}

% Un terme technique cité pour la première fois.
\newcommand{\terme}[1]{\emph{#1}}
% Un identifiant de code (classe, table, endpoint).
\newcommand{\code}[1]{\texttt{\small #1}}
